\chapter{Excessive Hunger (Polyphagia)}

\section*{Definition}
Polyphagia is pathologically increased appetite and food consumption beyond normal caloric needs. It represents a metabolic or endocrine disturbance in appetite regulation, most commonly associated with insulin deficiency or resistance.

\section*{What Happens in Your Body}
Appetite is regulated by the arcuate nucleus of the hypothalamus through orexigenic (neuropeptide Y, agouti-related peptide) and anorexigenic (proopiomelanocortin, cocaine-amphetamine-regulated transcript) pathways. Insulin and leptin suppress appetite; ghrelin stimulates it. In diabetes mellitus, cellular glucose deprivation despite hyperglycemia triggers counter-regulatory responses that stimulate hunger centers. Thyroid hormone excess increases metabolic rate, driving compensatory hyperphagia.

\section*{Causes}
\begin{itemize}
  \item Diabetes mellitus (type 1 and type 2 -- most common cause)
  \item Hyperthyroidism (thyrotoxicosis increases metabolic rate)
  \item Hypoglycemia (reactive, fasting, or medication-induced)
  \item Medications: Corticosteroids, antipsychotics (olanzapine, clozapine), antidepressants, antihistamines, cannabis
  \item Premenstrual syndrome (hormonal appetite fluctuations)
  \item Bulimia nervosa, binge eating disorder
  \item Prader-Willi syndrome (genetic appetite dysregulation)
  \item Hypercortisolism (Cushing syndrome)
  \item Sleep deprivation (disrupted leptin/ghrelin balance)
\end{itemize}

\section*{When to See a Doctor}
See your doctor if:
\begin{itemize}
  \item Polyphagia with unintentional weight loss (suggests diabetes or hyperthyroidism)
  \item Excessive thirst and urination accompanying increased appetite (diabetes mellitus)
  \item Heat intolerance, tremor, palpitations, or weight loss despite increased appetite (hyperthyroidism)
  \item Episodic extreme hunger with sweating, confusion, or shakiness (hypoglycemia)
  \item Binge eating episodes with loss of control (eating disorder)
\end{itemize}

\section*{Related Symptoms}
\begin{itemize}
  \item Weight loss
  \item Weight gain
  \item Fatigue
  \item Excessive thirst
  \item Frequent urination
  \item Tremor
  \item Heat intolerance
  \item Palpitations
\end{itemize}

\section*{Self-Care / Home Management}
\begin{itemize}
  \item Keep a food diary to document appetite patterns and triggers.
  \item Eat balanced meals with protein and fiber to promote satiety.
  \item Monitor capillary blood glucose if diabetic.
  \item Avoid high-glycemic-index foods that cause rapid glucose fluctuations.
  \item Establish regular sleep schedule to regulate appetite hormones.
  \item Seek medical evaluation for polyphagia with weight change, tremor, or excessive thirst.
\end{itemize}

\section*{How Your Doctor Will Evaluate You}
\begin{itemize}
  \item Fasting blood glucose, HbA1c, and oral glucose tolerance test to evaluate for diabetes.
  \item Thyroid function tests (TSH, free T4) to screen for hyperthyroidism.
  \item Cortisol levels (AM serum, 24-hour urine, dexamethasone suppression) if Cushing syndrome suspected.
  \item Nutritional assessment and eating disorder screening questionnaires.
\end{itemize}

\section*{Treatment Options}
\begin{itemize}
  \item Treat underlying endocrine disorder (insulin/oral hypoglycemics for diabetes, antithyroid drugs for hyperthyroidism).
  \item Nutritional counseling and structured meal planning.
  \item Cognitive behavioral therapy for eating disorders.
  \item Medication adjustment if drug-induced.
  \item Appetite-regulating medications in selected cases under specialist guidance.
\end{itemize}

\section*{References}
\begin{thebibliography}{99}
\bibitem{ref1} Schwartz MW, Seeley RJ, Zeltser LM, et al. "Obesity Pathogenesis: An Endocrine Society Scientific Statement." Endocr Rev. 2017;38(4):267--296.
\bibitem{ref2} Kliegman RM, St Geme JW III, Blum NJ, et al. "Nelson Textbook of Pediatrics." 22nd ed. Elsevier; 2024.
\bibitem{ref3} Melmed S, Auchus RJ, Goldfine AB, et al. "Williams Textbook of Endocrinology." 14th ed. Elsevier; 2020.
\bibitem{ref4} Blundell JE, Gillett A. "Control of Food Intake in the Obese." Obes Res. 2001;9(Suppl 4):263S--270S.
\bibitem{ref5} Saper CB, Lowell BB. "The Hypothalamus." Curr Biol. 2014;24(23):R1111--R1116.
\bibitem{ref6} Morton GJ, Meek TH, Schwartz MW. "Neurobiology of Food Intake in Health and Disease." Nat Rev Neurosci. 2014;15(6):367--378.
\end{thebibliography}
\clearpage
